\section{Source audit: qualifications, corrections, and non-theorem material}
\label{sec:source-audit}

The goal is not merely to reproduce formulas but to determine what can be proved with their stated hypotheses.  The following audit is part of the mathematical result.

\subsection{Corrections and required qualifications}

\begin{longtable}{@{}p{0.25\textwidth}p{0.66\textwidth}@{}}
\toprule
Source item & Constructive audit\\
\midrule\endhead
Wikipedia symmetric series for $\B(z,z)$ &
False as rendered with the explicitly stated rising-factorial convention.  At $z=1$ the first series has terms not tending to zero; the second Abel expression equals $2/3$, not $1$.  See \cref{warn:false-series}.\\

Fractional-linear integral with parameter $c\le1$ &
The displayed substitution requires $c<1$.  At $c=1$ it degenerates and the ordinary integrand/prefactor formula is undefined.  A limiting statement would need separate hypotheses.\\

Central radical formula with $z\ne1$ &
For positive real $z$ the proof works also at $z=1$; the exclusion is unnecessary.  For complex $z$, roots require branch choices that the source does not state.\\

Reflection $B(x,1-x)=\pi/\sin\pi x$ for $x\notin\Z$ &
The integral proof requires $0<\Rea x<1$.  Elsewhere it is meromorphic.  Constructive division requires a positive distance from the relevant integers, not merely negated membership.\\

Binomial interpolation quotient &
The raw quotient has $0/0$ at integer $k$.  Use the cleared entire identity or the factorial value at integers.\\

Incomplete-beta transformation with $(-1)^a$ &
Branch-ambiguous for noninteger complex $a$ and a negative transformed endpoint.  It is directly valid for integer $a$, or with an explicit branch/path certificate.\\

Continued-fraction bounds &
The finite continuant algebra is elementary.  The alternating inequalities and rapid-convergence statement require positivity and a remainder theorem; they are not consequences of formal continued-fraction notation alone.\\

MathWorld factorial notation &
Expressions such as $(p-1)!$ are literal factorials only at positive integers.  For general complex parameters they abbreviate Gamma values.\\

MathWorld polar derivation &
Correct classically, but unnecessary for the iPRA target.  A finite polar-grid proof needs a Jacobian error analysis; the one-dimensional beta-to-gamma limit is simpler.\\

Software and application paragraphs &
Statements about Excel, SciPy, string scattering, or preferential attachment are factual/application claims, not beta-function identities.  They are not formalized in this packet.\\
\bottomrule
\end{longtable}

\subsection{The corrected counterexample calculation}

For clarity, at $z=1$ the second displayed symmetric expression is
\[
 2\sum_{n\ge0}\frac{(-2)_n r^n}{(n+1)n!}.
\]
The rising factorial terminates:
\[
 (-2)_0=1,\qquad(-2)_1=-2,\qquad(-2)_2=2,\qquad(-2)_n=0\ (n\ge3).
\]
Therefore the $r\to1^-$ limit is
\[
 2\left(1-1+\frac13\right)=\frac23.
\]
The exact-rational evidence file checks this arithmetic.  This corrects an earlier scratch calculation and is the value used throughout the final packet.

\subsection{Statements not fully proved here}

The following are mathematically plausible standard results but remain at status \statusC{} in this packet:

\begin{enumerate}
\item the Pochhammer-contour identity with explicit branch and polygonal-contour certificates;
\item the reciprocal-beta trigonometric integrals for arbitrary noninteger frequency;
\item the general alternating-convergent remainder theorem for the incomplete-beta continued fraction;
\item a completely internalized variable-dimension simplex integration theorem;
\item a kernel-checked production of all analytic statements as \code{Prf PRAAx} terms.
\end{enumerate}

These are not omitted because of classical logic.  Each has a finitary route; the missing work is the construction and checking of substantial quantitative certificates.

\subsection{Why the audit matters for iPRA}

Classical expositions often use analytic continuation to make a formula true after silently changing its meaning.  That practice is especially dangerous in a weak arithmetic setting.  An integral, a meromorphic quotient, a branch-dependent contour value, and a removable-singularity extension are four different algorithms.  They may agree on overlapping domains, but the agreement itself needs a theorem and a modulus.  This report therefore states domains at every transition and favors cleared equations at possible singularities.
